\section{数值结果}
\label{sec:num}

本节我们通过研究式~(\ref{eq:fac-tq},\ref{eq:ps-q-fact})中的匹配系数如何改变 PDF，对准 PDF、\spcorr 和赝 PDF 进行数值分析。对 $\overline{\rm MS}$、横向动量截断与 RI/MOM 三种方案的准 PDF，文献~\cite{Stewart:2017tvs}已做过这类分析。我们新的 $\overline{\rm MS}$ 匹配结果由式~(\ref{eq:quasi-matching})给出，它带来稳定的卷积积分。我们还比较了 $\overline{\rm MS}$ 方案中匹配系数使用强子动量 $p^z=P^z$ 与部分子动量 $p^z=|y|P^z$ 的差异。这里的结果取 $\Gamma =\gamma^0$。

作为分析示例，我们采用非极化同位旋矢量部分子分布
\begin{align}
f_{u-d}(x,\mu) =  f_u(x,\mu) - f_d(x,\mu) - f_{\bar{u}}(-x,\mu) +  f_{\bar{d}}(-x,\mu),
\end{align}
其中包含反部分子分布 $f_{\bar{u}}(-x,\mu) = - f_{\bar{u}}(x,\mu)$ 与 $f_{\bar{d}}(-x,\mu) = - f_{\bar{d}}(x,\mu)$。为便于比较，我们使用 MSTW 2008~\cite{Martin:2009iq} 的次领头阶同位旋矢量 PDF $f_{u-d}$ 及相应的跑动耦合 $\alpha_s(\mu)$。

为了在数值计算中实现 plus 函数，我们施加软截断 $|y-x|<10^{-m}$ 并检验结果对 $m$ 的敏感性。由于 $y\to0$ 对应渐近区 $|x/y|\to\infty$，我们还施加紫外截断 $|y|>10^{-n}$ 来检验卷积积分的收敛性。我们发现下面给出的所有结果对 $m$ 与 $n$ 都不敏感。我们的式~(\ref{eq:quasi-matching})的结果在 $\xi\in [1,\infty]$ 与 $\xi\in [-\infty,0]$ 两个区间内都有位于 plus 函数之外的项，这对保证我们的 $\overline{\rm MS}$ 结果 $C$ 对 $|y|>10^{-n}$ 截断不敏感十分重要。对于以横向动量截断定义的准 PDF 则不然~\cite{Xiong:2013bka}。RI/MOM 方案的结果也不受此问题困扰~\cite{Stewart:2017tvs}。

在图~\ref{fig:quasi}中，我们把 PDF 与由式~(\ref{eq:fac-tq})的卷积、用我们的单圈结果式~(\ref{eq:quasi-matching})得到的 $\overline{\rm MS}$ 准 PDF 作了比较。我们观察到，从 $p^z=P^z$ 改为正确的 $p^z=|y|P^z$ 会使物理区域内的结果发生相当大的移动。

\begin{figure*}
	\centering
	\includegraphics[width=.49\textwidth]{images/Compare_iPDFvsPDFRe_TeX.pdf}
	\hspace{0.1cm}
	\includegraphics[width=0.49\textwidth]{images/Compare_iPDFvsPDFRe_z_TeX.pdf}
	\\
	\includegraphics[width=.49\textwidth]{images/Compare_iPDFvsPDFIm_TeX.pdf}
	\hspace{0.1cm}
	\includegraphics[width=0.49\textwidth]{images/Compare_iPDFvsPDFIm_z_TeX.pdf}	\\
\vspace{-0.3cm}
	\caption{（左）光锥时间分布 $Q_{u-d}$ 与由 $({\cal C}^{\overline{\rm MS}}\otimes Q_{u-d})$ 得到的 $\overline{\rm MS}$ 方案\spcorr 的比较。橙色带与蓝色带表示因子化标度 $\mu=4\,{\rm GeV}$ 变化一倍的结果。（右）同左图，但只显示不同 $z$ 值的中心\spcorr 曲线。上排为实部，下排为虚部。 }
	\label{fig:ioffe}
\end{figure*}

同样的比较也可以对 $\overline{\rm MS}$ 方案的赝 PDF 进行：应用因子化公式式~(\ref{eq:ps-q-fact})与匹配系数式~(\ref{eq:ps-c})。在图~\ref{fig:pseudo}a 中我们比较 PDF 与赝 PDF 及其对因子化标度 $\mu$ 的依赖，而在图~\ref{fig:pseudo}b 中加入赝 PDF 对距离 $|z|$ 的依赖。由于式~(\ref{eq:ps-c})中的匹配系数除了非平凡的有限常数外类似于部分子劈裂函数，把 PDF 匹配到赝 PDF 类似于把 PDF 从 $\mu$ 演化到 $1/|z|$ 标度。这一演化已在文献~\cite{Radyushkin:2017cyf}中计算。如图~\ref{fig:pseudo}右图所示，$|z|$ 的变化具有与 PDF 演化类似的效应。当 $|z|\mu=1$ 时对数为零，${\cal C}$ 的匹配效应由式~(\ref{eq:ps-c})中的非平凡常数决定，它使 PDF 在大 $x$ 区域下移、在小 $x$ 区域上移。

最后，我们可以对用式~(\ref{eq:io-Q-fact})与式~(\ref{eq:ps-c})得到的 $\overline{\rm MS}$ 方案\spcorr 作类似比较。其实部与虚部分别是 $\zeta$ 的偶函数与奇函数，示于图~\ref{fig:ioffe}。我们同样给出对 $\mu$ 与 $|z|$ 的残余依赖，它与赝 PDF 的情形相似。匹配展宽了坐标空间中的曲线。$\overline{\rm MS}$ 方案重正化的\spcorr 不表现出矢量流（或粒子数）守恒，这可以清楚地从分布在 $\zeta=0$ 处实部不为 1 看出（除非 $|z|\mu$ 恰好调到抵消单圈 ${\cal C}$ 中常数项的特殊情形）。
